\section{Methodology}
\label{sec:method}

\subsection{Semantic Segmentation}
\begin{table}[!htbp]
\centering
\small
\caption{Semantic segmentation results on \texttt{Cityscapes val} (mIoU \%). 
EoMT-COCO is evaluated zero-shot with COCO$\rightarrow$Cityscapes class mapping. 
EoMT-Fine-tuned starts from EoMT-COCO and is fine-tuned on Cityscapes.}
\label{tab:semseg}
\begin{tabular}{lc}
\toprule
Model & mIoU \\
\midrule
ERFNet                  & 72.50 \\
EoMT-COCO (zero-shot)   & 55.17 \\
EoMT-Cityscapes         & \textbf{81.68} \\
EoMT-Fine-tuned         & 78.85 \\
\bottomrule
\end{tabular}
\end{table}

Semantic segmentation was our first task, and we used EoMT~\cite{eomt}, which is a mask-based architecture built on a DINOv2-pretrained Vision Transformer~\cite{dinov2}. 
We compared two EoMT models using two different checkpoints: EoMT-COCO and EoMT-Cityscapes. EoMT-COCO is trained on the COCO dataset~\cite{coco}, a large general-purpose dataset with 118k training images covering daily scenes such as streets, kitchens, beaches, and animals. 
COCO has 133 class labels: 80 ``thing'' classes such as person, car, dog, and 53 ``stuff'' classes such as sky, grass, wall. 
This makes the model strong on a wide range of visual objects, but it is not road-specific.
The second model is EoMT-Cityscapes, trained on the Cityscapes dataset~\cite{cityscapes}. Cityscapes is a smaller dataset (around 3k training images) of road and driving scenes from German cities. 
It has 19 class labels specifically chosen for autonomous driving, such as road,traffic sign, rider. It is road-specific but not as broad as COCO.
We evaluated both models on the Cityscapes validation set, which contains 500 images. 
The goal is to understand how much a general model loses compared to a specialized one, and whether fine-tuning close the gap. 
We need to translate the COCO predictions into the Cityscapes label space before evaluation.
We built a \texttt{coco\_to\_cityscapes} mapping that works in three steps:
\begin{itemize}
    \item \textbf{Direct name match:} obvious matches found by string comparison. For example, COCO ``car'' maps directly to Cityscapes ``car'', and COCO ``person'' to Cityscapes ``person''.
    \item \textbf{Synonyms:} some classes are named differently in the two datasets. We built a synonym dictionary to handle these cases. For example, COCO includes sidewalks in ``pavement-merged''.
    \item \textbf{Ignored:} COCO labels with no Cityscapes equivalent are mapped to the ignore index 255, which excludes them from evaluation. For example, if EoMT-COCO predicts ``dog'', that pixel is ignored.
\end{itemize}
Intersection-over-Union (IoU) is a method to measure the accuracy for each class:
\begin{equation}
\mathrm{IoU}_{c} = \frac{TP_{c}}{TP_{c} + FP_{c} + FN_{c}}
\label{eq:iou}
\end{equation}
Mean IoU (mIoU) is the average of these scores.

\subsection{Fine-Tuning} 
We wanted to see if we could adapt it to the Cityscapes domain without losing its general knowledge.
For fine tuning, we drop the classification head because COCO's classification head produces 133 logits per query. We replaced it with a new head that has 19 outputs which is one per Cityscapes class.
The rest of the weights are kept frozen.
For efficiency, we froze the DINOv2 backbone and trained the queries, and the new classification head. 
We also use Automatic Mixed Precision (AMP). Instead of only FP32, we use mixed-precision.
For the training setup $20$ epochs, $K=200$ queries, image size $640{\times}640$ used. 
Training was done on a single NVIDIA L4 GPU via Google Colab Pro.
For the loss function we used standard Mask2Former~\cite{mask2former} loss.
We use AdamW with a learning rate of $1\times10^{-4}$.

\subsection{Anomaly Segmentation}
We evaluate ERFNet and EoMT under anomaly segmentation.
Anomaly segmentation detects pixels belonging to objects which are not found in the training class set, out-of-distribution.
For evaluation, we use post-hoc methods which use class logits, and calculate anomaly score per pixel. 
Both ERFNet (pixel-based) and EoMT (mask-based) produce class logits, and this allows us to use the same scoring functions for both models.
\paragraph{Post-hoc anomaly detection methods.}
\begin{itemize}
    \item \textbf{MSP (Maximum Softmax Probability):} uses the highest softmax probability as an uncertainty indicator.
    \begin{equation}
    \text{score}_{\text{MSP}}(\mathbf{x}) = 1 - \max_c p_c(\mathbf{x})
    \end{equation}
    
    \item \textbf{MaxLogit}~\cite{maxlogit}: uses the raw logits before softmax, because they preserve the magnitude of the model's confidence.
    \begin{equation}
    \text{score}_{\text{MaxLogit}}(\mathbf{x}) = -\max_c L_c(\mathbf{x})
    \end{equation}
    
    \item \textbf{MaxEntropy}~\cite{maxentropy}: if the probability is spread across many classes (high entropy), the model is uncertain and the pixel is a possible anomaly. If it puts all probability on one class (low entropy), the model is confident.
    \begin{equation}
    \text{score}_{\text{Entropy}}(\mathbf{x}) = -\sum_{c=1}^{C} p_c(\mathbf{x}) \log p_c(\mathbf{x})
    \end{equation}
    
 \item \textbf{\textcolor{blue}{RbA (Rejected by All)~\cite{rba}:}} \textcolor{blue}{RbA is used only for mask-based models. It is based on the idea that an   unknown pixel should receive low acceptance from all known classes. In our pipeline, the EoMT per-class semantic score map is already obtained by combining mask predictions and class predictions while excluding the no-object class. Therefore, in the revised implementation, we compute RbA as the negative total known-class mass.}
\textcolor{blue}{In a pixel-wise classifier, class probabilities are normalized by the softmax and must sum to 1, so the rejected-by-all interpretation is not valid in the same way. That is why we do not apply RbA to ERFNet.}
\begin{equation}
\textcolor{blue}{
\text{score}_{\text{RbA}}(\mathbf{x}) = -\sum_{c=1}^{C} L_c(\mathbf{x})
}
\end{equation}
\textcolor{blue}{Here, \(L_c(\mathbf{x})\) is the EoMT semantic score for class \(c\) at pixel \(\mathbf{x}\). Higher scores indicate lower total acceptance by known classes and therefore higher anomaly likelihood.}
\end{itemize}


\paragraph{Temperature scaling.}

It is a calibration technique that can shift recall and precision in anomaly detection.
It divides logits by $T$ before softmax. $T > 1$ makes probabilities more uncertain, $T < 1$ makes them more confident. 
\textcolor{blue}{In the revised implementation, we tested $T \in \{0.5, 0.75, 1.0, 1.1, 1.5, 2.0\}$ only on MSP and MaxEntropy. We did not test temperature scaling on MaxLogit because it doesn't use softmax, and we didn't apply it to RbA because the revised RbA score is based on the total known-class mass rather than softmax calibration.}

\paragraph{Datasets and metrics.}

As evaluation benchmarks, we use a smaller subsets of these five datasets which are provided by course instructors. These are:

\begin{itemize}
    
    \item \textbf{RoadAnomaly21 (from SMIYC~\cite{smiyc}):} 10 images, large unknown objects in diverse environments.
    \item \textbf{RoadObstacle21 (from SMIYC~\cite{smiyc}):} 30 images, small unknown objects on the road.
    \item \textbf{Fishyscapes Lost\&Found~\cite{fishyscapes}:} 100 validation images of real lost-cargo objects.
    \item \textbf{Fishyscapes Static~\cite{fishyscapes}:} 30 images, synthetic anomalies on the road scenes.
    \item \textbf{RoadAnomaly (from SMIYC~\cite{smiyc}):} 60 images, earlier version of SMIYC's anomaly track.     

\end{itemize}
We use two metrics:
\begin{itemize}
    \item \textbf{AuPRC} (Area under the Precision-Recall Curve): Evaluates pixel-level separability. Higher values mean better, used for discrimination of rare anomalies.
    \item \textbf{FPR@95} (False Positive Rate at 95\% True Positive Rate): Measures false positive rate when the true positive rate is fixed at 95\%. It is a safety critical metric. Lower is better.
\end{itemize}
